\subsection*{Expression Topography: Where Risk Genes Meet Regulatory Circuits}

The genome-wide association catalog and candidate gene literature converge on a
handful of loci implicated in emotional dysregulation, yet a gene's name on a
Manhattan plot tells us almost nothing about \emph{where} it acts. Two risk
genes may sit adjacent on a chromosome and operate in entirely different neural
circuits; conversely, genes on separate chromosomes may converge on the same
synapse. To move from statistical association to mechanistic explanation, we
must ask a deceptively simple question: in which cells, in which brain regions,
is each risk gene actually expressed? The answer, as we shall see, reorganizes
the prevailing circuit model of emotional dysregulation in ways that no amount
of functional MRI alone could achieve.

Consider \gene{SLC6A4}, encoding the serotonin transporter, arguably the most
studied gene in psychiatric genetics. The standard narrative places it at the
center of the prefrontal--amygdala circuit: variation in \gene{SLC6A4}
modulates amygdala reactivity to threat \cite{hariri2002serotonin}, alters
prefrontal--amygdala functional connectivity
\cite{pezawas2005slc6a4connectivity}, and interacts with early adversity to
predict depressive episodes \cite{caspi2003influence}. One might reasonably
assume that the transporter is expressed in amygdala neurons, clearing serotonin
from local synapses. The expression data say otherwise.

Queries to the Human Protein Atlas for \gene{SLC6A4} messenger RNA in the
amygdala return no detected expression \cite{uhlen2015tissue}. The same null
result emerges from the Genotype-Tissue Expression (GTEx) consortium: across
thirteen brain subregions profiled by GTEx, \gene{SLC6A4}
(ENSG00000108576.14) yields effectively empty returns
\cite{gtex2020gtex}. These are not failures of detection. They are informative
absences. The Allen Human Brain Atlas and its mouse homolog confirm that
\gene{Slc6a4} expression is overwhelmingly concentrated in the dorsal and
median raphe nuclei of the brainstem, with modest expression in select cortical
and limbic territories in mouse \cite{lein2007genome, hawrylycz2012anatomically}.
The raphe nuclei---small, phylogenetically ancient cell clusters---are the
principal source of serotonergic innervation throughout the forebrain. They are
also, notably, absent from the standard GTEx tissue panel, which explains the
consortium-wide silence.

The mechanistic implication is fundamental. \gene{SLC6A4} does not modulate
amygdala function from within the amygdala. It does so \emph{remotely}, from
serotonergic cell bodies in the raphe nuclei whose axonal projections terminate
in the amygdala, the medial prefrontal cortex, and the hippocampus
\cite{hale2012functional}. The transporter protein is synthesized in the raphe
soma, trafficked down the axon, and inserted into presynaptic membranes at
distant terminal fields, where it clears serotonin from the synaptic cleft.
Genetic variation that alters transcription, splicing, or protein folding of
\gene{SLC6A4} therefore acts at the \emph{source}---the raphe---but manifests
functionally at the \emph{target}: the amygdala terminal, the prefrontal
terminal, and every other serotonergic synapse in the brain simultaneously.

This three-node architecture---raphe to prefrontal cortex, raphe to
amygdala---reframes how we understand the transdiagnostic reach of
\gene{SLC6A4} variation. The prefrontal--amygdala model of emotional
dysregulation typically emphasizes a single dyadic relationship: top-down
prefrontal regulation of amygdala reactivity. But if the serotonin transporter
gene is expressed in neither node of that dyad, then its influence must be
understood as a simultaneous modulation of \emph{both} arms of the circuit from
a common upstream source. A low-expressing \gene{SLC6A4} allele does not
selectively disinhibit the amygdala; it alters serotonergic tone at every
terminal field the raphe innervates, including prefrontal regions responsible
for reappraisal, hippocampal circuits supporting contextual fear conditioning,
and striatal pathways mediating reward sensitivity
\cite{murphy2008serotonin, canli2006long}. This \emph{broadcast}
architecture---one gene, one cell group, many targets---explains why
\gene{SLC6A4} polymorphisms surface in studies of depression, anxiety,
impulsivity, aggression, and borderline personality features with roughly
comparable, if modest, effect sizes \cite{karg2011serotonin}. The gene is not
pleiotropic in the classical sense of acting through multiple molecular
pathways; it is pleiotropic by \emph{projection geometry}.

The selective serotonin reuptake inhibitors (SSRIs) exploit exactly this
geometry, though the pharmacological framing rarely makes the point explicit.
The drug acts at the synaptic terminal---blocking the transporter protein at
the amygdala synapse, the prefrontal synapse, the hippocampal synapse---yet the
protein it targets is encoded by a gene expressed in the raphe. The therapeutic
and the genetic thus operate at opposite ends of the same axon. This has
practical consequences for pharmacogenomics: a functional polymorphism that
reduces \gene{SLC6A4} transcription in raphe neurons will lower transporter
density at \emph{all} terminal fields, potentially altering SSRI binding
kinetics across the entire serotonergic projection system rather than at any
single circuit node \cite{murphy2004pharmacogenomics}.

The contrast with the glucocorticoid-signaling genes is instructive.
\gene{NR3C1}, encoding the glucocorticoid receptor, and \gene{FKBP5}, encoding
its cochaperone FK506-binding protein 51, show robust expression in the
hippocampus and prefrontal cortex---the very structures whose function they are
proposed to modulate \cite{perlman2007glucocorticoid, binder2009association}.
Variation in these genes alters the stress response \emph{locally}, within the
neurons that integrate hypothalamic--pituitary--adrenal (HPA) axis feedback.
Hippocampal neurons expressing \gene{NR3C1} bind circulating cortisol and
adjust their own firing, synaptic plasticity, and gene expression in direct
response \cite{deangelis2019glucocorticoids}. There is no obligatory relay
through a distant cell group. The gene and the circuit are coextensive.

These two expression architectures---distal projection modulation by
\gene{SLC6A4} versus local circuit action by \gene{NR3C1} and
\gene{FKBP5}---are not merely neuroanatomical curiosities. They generate
distinct predictions for gene--environment interaction, pharmacological
response, and the developmental timing of risk. A gene expressed at the circuit
node can be regulated by local epigenetic signals: childhood adversity alters
\gene{NR3C1} methylation in hippocampal neurons, directly silencing the
receptor in the cells that need it \cite{mcgowan2009epigenetic}. A gene
expressed at a distal projection source, by contrast, is subject to the
epigenetic environment of the \emph{raphe}, a region whose developmental
programming and chromatin landscape remain poorly characterized in humans.
Until brainstem-resolution epigenomic atlases become available, the
environmental regulation of \gene{SLC6A4} will remain substantially opaque.

The expression topography of risk genes thus reveals a deeper structure beneath
the familiar circuit diagrams. Emotional dysregulation is not simply a matter
of weakened prefrontal control over amygdala reactivity. It is a network
phenomenon shaped by genes that act at different spatial scales: some from
within the circuit, tuning local synapses; others from distant brainstem nuclei,
simultaneously adjusting the neurochemical milieu across multiple forebrain
targets. Integrating transcriptomic atlases with connectomic and epigenomic
data will be essential for building genetic models of emotional dysregulation
that honor this spatial complexity---models in which the locus of expression,
not merely the locus on the chromosome, determines the architecture of risk.
